\section{Learning Channel and Activity Denoisers}
\label{sec:learning}

Training first learns an AWGN channel score and then adapts independent
layer instances using queries produced by the unfolded receiver.

\subsection{First-Order Score Matching}

Let $h_R\in\mathbb R^{2D}$ be a normalized channel before activity
masking. For $c=v/2$, generate $u=h_R+\sqrt c\,\varepsilon$ with
$\varepsilon\sim\mathcal N(0,\I_{2D})$. If $q_{1,R}(u;v)$ is its
density, the score is $s^\star(u,v)=\nabla_u\log q_{1,R}(u;v)$.
Differentiating the Gaussian convolution gives
\begin{equation}
 \nabla_u q_{1,R}(u;v)
 =q_{1,R}(u;v)\frac{\E[h_R\mid u]-u}{c},\qquad
 \E[h_R\mid u]=u+c\,s^\star(u,v).
 \label{eq:twc-tweedie}
\end{equation}
This is the real-coordinate form of Tweedie's formula
\cite{efron2011,cai2026}. Denoising score matching
\cite{vincent2011,song2019} trains the network by
\begin{equation}
 \mathcal L_{\rm DSM}(\theta)=\E_{h_R,v,\varepsilon}
 \left[\frac{\|\sqrt{v/2}\,s_\theta(h_R+\sqrt{v/2}\varepsilon,v)
                      +\varepsilon\|^2}{2D}\right].
 \label{eq:twc-dsm}
\end{equation}
The distribution of $v$ is fixed before training. The conditional
mean of the regression target is the score, so this objective recovers
$s^\star$ in the unrestricted function class on its training support.
Only active-channel samples enter prior training; inactive zeros are
represented by the separate activity model.

Equation~\eqref{eq:twc-tweedie} explains the initialization of
\eqref{eq:twc-conditional-mean} at $a_\ell=1$. After receiver
fine-tuning, the network is a task-adapted denoiser in this
parameterization and need not remain the exact score of the AWGN prior.

\subsection{Conditional Channel and Activity Supervision}

The receiver training set has $B$ scenes and $K_i$ users in scene $i$.
Define fixed training statistics
\begin{equation}
 T_{\rm tr}=\sum_iK_i,\qquad
 E_{\rm tr}=\sum_{i,k}\alpha_{ik}\|h_{ik,R}\|^2,\qquad
 \bar K=T_{\rm tr}/B,\quad \bar E=E_{\rm tr}/B.
 \label{eq:twc-training-statistics}
\end{equation}
Here $E_{\rm tr}>0$, and the real-coordinate squared norm equals
complex channel energy. For each layer use
\begin{align}
 R_H^{(\ell)}&=\frac{1}{E_{\rm tr}}\sum_{i,k}
     \alpha_{ik}\|\mu_{ik\ell}-h_{ik,R}\|^2,
 \label{eq:twc-channel-loss}\\
 R_a^{(\ell)}&=\frac{1}{T_{\rm tr}}\sum_{i,k}
     {\rm BCE}(\alpha_{ik},p_{ik\ell}),
 \label{eq:twc-activity-loss}\\
 {\rm BCE}(\alpha,p)&=-\alpha\log p-(1-\alpha)\log(1-p).
 \label{eq:twc-bce}
\end{align}
True activity selects conditional channel supervision only during
training. All users enter unweighted BCE; inference forms
$\widehat x=p\mu$ using the predicted probability.
Effective channel error is evaluated separately.

Fix the information $z$ available to a user estimator and let
$\pi(z)=P(\alpha=1\mid z)$,
$m_H(z)=\E[h_R\mid z,\alpha=1]$, and
$\kappa=E_{\rm tr}/T_{\rm tr}>0$.
For finite conditional second moments, $0<\pi(z)<1$, and arbitrary
outputs $\mu$ and $p\in(0,1)$,
\begin{equation}
 \begin{split}
 &\E\left[\frac{\alpha}{\kappa}\|\mu-h_R\|^2+
                       {\rm BCE}(\alpha,p)\mid z\right]\\
 &\quad=C(z)+\frac{\pi(z)}{\kappa}\|\mu-m_H(z)\|^2+
 {\rm KL}({\rm Bern}(\pi(z))\,\|\,{\rm Bern}(p)).
 \end{split}
 \label{eq:twc-risk-decomposition}
\end{equation}
Here $C(z)$ is independent of the outputs and $\operatorname{KL}$ is
Kullback--Leibler divergence. Appendix~\ref{app:risk} gives the
expansion. The unrestricted minimizers are
\begin{equation}
 p^\star(z)=\pi(z),\qquad \mu^\star(z)=m_H(z),\qquad
 p^\star(z)\mu^\star(z)=\E[\alpha h_R\mid z].
 \label{eq:twc-risk-minimizers}
\end{equation}
This explains the effective channel target without adding a third
loss. It does not establish optimization or generalization of finite
networks, where outputs share dependencies and queries change with
earlier layers. Fixed $E_{\rm tr}$ is essential: division by a
sample-dependent true channel energy generally changes the target.

For uniform scene sampling, implement
\begin{equation}
 \ell_{H,i}^{(\ell)}
 =\frac{\sum_k\alpha_{ik}\|\mu_{ik\ell}-h_{ik,R}\|^2}{\bar E},\qquad
 \ell_{a,i}^{(\ell)}
 =\frac{\sum_k{\rm BCE}(\alpha_{ik},p_{ik\ell})}{\bar K}.
 \label{eq:twc-scene-loss}
\end{equation}
Their scene averages recover
\eqref{eq:twc-channel-loss}--\eqref{eq:twc-activity-loss}, including
when $K_i$ varies. Validation uses the same training constants.

\subsection{Layer Training and Gradient Propagation}

At depth $m$, minimize
\begin{equation}
 J_m=\frac{\sum_{\ell=1}^{m}\omega_\ell
                     (R_H^{(\ell)}+R_a^{(\ell)})}
                 {\sum_{\ell=1}^{m}\omega_\ell}.
 \label{eq:twc-layer-objective}
\end{equation}
Equal coefficients after normalization specify a task weighting,
not a uniquely optimal choice. Layer weights are fixed before training.
Each new layer independently loads the same pretrained score weights.
Earlier layers retain the weights selected at the preceding depth.
Classifier output corrections and log calibrations are initialized
to zero, giving $p_{k\ell}=\rho$. Each depth first adapts the new
layer and then jointly fine-tunes all layers at a lower learning rate.
The progressive training principle follows \cite[Sec.~III-D]{wei2026};
conditional channel supervision, activity BCE, and the stage schedule
are specified for the present JADCE receiver.

Let $\Theta_m$ contain all score, activity-head, and calibration
parameters that influence the depth-$m$ outputs. Training uses
\begin{equation}
 \nabla_{\Theta_m}J_m
 =\nabla_{\Theta_m}J_{H,m}+\nabla_{\Theta_m}J_{a,m},
 \label{eq:twc-joint-gradient}
\end{equation}
where $J_{H,m}$ and $J_{a,m}$ are the two weighted loss terms.
BCE gradients can reach same-layer and earlier score parameters.
Conditional channel loss does not depend on the same-layer activity
head, but later channel losses can update earlier heads through the
messages. Freezing parameters during adaptation does not detach
their input derivatives.

The final layer's $c_{\beta,m}$ affects only feedback after its output,
so it is not optimized until another layer is appended.
New-layer adaptation also enables the preceding final layer's feedback
calibration. At each depth select the checkpoint minimizing
\begin{equation}
 V_m=\frac{1}{B_{\rm val}}\sum_{i\in{\rm val}}
     (\ell_{H,i}^{(m)}+\ell_{a,i}^{(m)}),
 \label{eq:twc-selection}
\end{equation}
The earlier checkpoint is chosen in a tie and initializes the next depth.
All primary final-depth metrics come from the same selected checkpoint;
the rule does not imply simultaneous optimality of detection and
estimation metrics.
